\pdfoutput=1
\documentclass{article}

%%% ready for submission %%%
% \usepackage{cpal_2025}

%%% to compile a preprint version, e.g., for submission to arXiv, add the [preprint] option %%%
% \usepackage[preprint]{cpal_2025}

%%% to compile a camera-ready version, add the [final] option %%%
\usepackage[final]{CPAL-2025-template/cpal_2025}


%add packages
\usepackage{hyperref}
\usepackage{url}
\usepackage{xspace}

\usepackage{booktabs}       % professional-quality tables
\usepackage{amsfonts}       % blackboard math symbols
\usepackage{nicefrac}       % compact symbols for 1/2, etc.
\usepackage{microtype}      % microtypography
\usepackage{xcolor}         % colors
\usepackage{xspace}
\usepackage{natbib}

\usepackage{graphicx}
\usepackage{pifont}
\usepackage{tcolorbox}
\usepackage{tabularx}
\usepackage{makecell}
\usepackage{multirow}
\usepackage{multicol}
\usepackage{amssymb}
\usepackage{bbm}
\usepackage{amsmath}
\usepackage{wrapfig}
\usepackage{cleveref}
\usepackage{enumitem}
\usepackage{titletoc}

\usepackage{mathtools}
\usepackage{amsthm}
\newtheorem{lemma}{Lemma}
\usepackage{array}
\usepackage{caption}
\usepackage{algorithm}
\usepackage{algorithmic}
\usepackage{graphicx}
\usepackage{booktabs}
\usepackage{xcolor}
\usepackage{multirow}
\usepackage{amsfonts}
\usepackage{xcolor}
\usepackage{colortbl}
\usepackage{colortbl}
\usepackage{makecell}
\usepackage{float}
\usepackage{pifont}
\usepackage{placeins}
\usepackage{wrapfig}
\usepackage{tabularx}
\usepackage{subcaption}
\usepackage[english]{babel}      % 开启英文断词规则
\usepackage{array}               % 支持自定义列格式
\usepackage[table]{xcolor}
\newcolumntype{R}[1]{>{\raggedright\arraybackslash}p{#1}} % 定义带自动换行的左对齐列
% \usepackage{bibunits}
\definecolor{lightergreen}{RGB}{228, 108, 10}


\definecolor{darkgreen}{rgb}{0.0, 0.5, 0.0}
\definecolor{text_red}{RGB}{220,20,60}

\hypersetup{
  colorlinks,
  linkcolor={red!90!black},
  citecolor={orange!90!black},
  urlcolor={blue!70!black}
}

\title{ERC-SVD: Error-Controlled SVD for Large Language Model Compression}


\author{%
    Haolei Bai\textsuperscript{1,2},
    ~Siyong Jian\textsuperscript{1,4},
    ~Tuo Liang\textsuperscript{3},
    ~Yu Yin\textsuperscript{3},
    ~Huan Wang\textsuperscript{1}\thanks{Corresponding author: \texttt{wanghuan@westlake.edu.cn}}\\
  \textsuperscript{1}Westlake University, \textsuperscript{2}Nanyang Technological University, \\ \textsuperscript{3}Case Western Reserve University, \textsuperscript{4}Nanjing University\\
  % \texttt{zhu.3440@osu.edu, cyou@google.com}
}

\begin{document}


\maketitle


\input{CPAL-2025-template/sec/0_abstract}
\input{CPAL-2025-template/sec/1_introduction}
\input{CPAL-2025-template/sec/2_related_work}
\input{CPAL-2025-template/sec/3_method}
\input{CPAL-2025-template/sec/4_experiment}
\input{CPAL-2025-template/sec/5_conclusion}

\section*{Acknowledgement}
This paper is supported by Young Scientists Fund of the National Natural Science Foundation of China (NSFC) (No. 62506305), Zhejiang Leading Innovative and Entrepreneur Team Introduction Program (No. 2024R01007), Key Research and Development Program of Zhejiang Province (No. 2025C01026), Scientific Research Project of Westlake University (No. WU2025WF003), Chinese Association for Artificial Intelligence (CAAI) \& Ant Group Research Fund - AGI Track (No. 2025CAAI-ANT-13). It is also supported by the research funds of the National Talent Program and Hangzhou Municipal Talent Program.

% Reference
% For natbib users:

% \clearpage
\phantomsection
\bibliography{CPAL-2025-template/reference}



%%%%%%%%%%%%%%%%%%%%%%%%%%%%%%%%%%%%%%%%%%%%%%%%%%%%%%%%%%%%
\clearpage
\phantomsection
\input{CPAL-2025-template/sec/6_appendix}

\end{document}
